% !TEX root = ../algebra_02.tex

\section{Теория полей. Гомоморфизмы колец}

\begin{Remark}{}{}
	Все кольца в данном конспекте предполагаются коммутативными, ассоциативными и с единицей.
\end{Remark}

\subsection{Подкольца и идеалы}

\begin{Definition}{Подкольцо}{}
	Пусть $R$ --- кольцо. Подмножество $S \subset R$ называется подкольцом, если:
	\begin{enumerate}
		\item $S$ является подгруппой по сложению в $R$;
		\item $\forall a, b \in S \colon ab \in S$ (замкнутость относительно умножения);
		\item $1_R \in S$, откуда следует, что $1_S = 1_R$.
	\end{enumerate}
\end{Definition}

\begin{Example}{}{}
	Рассмотрим кольцо $R = \mathbb{Z} \times \mathbb{Z}$ и подмножество $S = \mathbb{Z} \times \{0\}$. 
	Единица кольца $R$ равна $1_R = (1, 1)$, однако в $S$ нейтральный элемент по умножению равен $1_S = (1, 0)$. Поскольку $1_R \notin S$, подмножество $S$ не является подкольцом в $R$ согласно нашему определению.
\end{Example}

\begin{Definition}{Идеал}{}
	Подмножество $I \subset R$ называется идеалом, если:
	\begin{enumerate}
		\item $I$ является подгруппой по сложению в $R$;
		\item $\forall a \in I, \forall x \in R \colon xa \in I$.
	\end{enumerate}
\end{Definition}

\begin{Remark}{}{}
	Если $K$ --- поле, то в $K$ существует ровно два идеала: $\{0\}$ и $K$. 
	Действительно, если $I \neq \{0\}$, то существует $a \in I \setminus \{0\}$. Так как $K$ --- поле, существует $a^{-1} \in K$. По определению идеала $1 = a^{-1}a \in I$, откуда $I = K$.
\end{Remark}

\subsection{Факторкольцо}

Пусть $I$ --- идеал в кольце $R$. Рассмотрим факторгруппу $R/I$ по сложению. Введем операцию умножения на классах смежности $\bar{a} = a + I$:
\[ \bar{a} \cdot \bar{b} = \overline{ab} \]

Покажем корректность этой операции. Пусть $\bar{a} = \bar{a}'$ и $\bar{b} = \bar{b}'$. Тогда $a' = a + x$ и $b' = b + y$ для некоторых $x, y \in I$. Имеем:
\[ a'b' = (a + x)(b + y) = ab + xb + ay + xy \]
Так как $I$ --- идеал, элементы $xb, ay, xy$ принадлежат $I$. Следовательно:
\[ \overline{a'b'} = \overline{ab} \]

\begin{Proposition}{}{}
	Множество классов смежности $R/I$ с введенными операциями сложения и умножения является кольцом.
\end{Proposition}

\begin{proof}
	Доказательство сводится к рутинной проверке аксиом кольца.
\end{proof}

\begin{Example}{}{}
	Факторкольцо кольца целых чисел по идеалу $n\mathbb{Z}$: $\mathbb{Z}/n\mathbb{Z} = \mathbb{Z}/(n)$.
\end{Example}

\subsection{Гомоморфизмы колец}

\begin{Definition}{Гомоморфизм колец}{}
	Пусть $R$ и $S$ --- кольца. Отображение $\varphi\colon R \to S$ называется гомоморфизмом колец, если:
	\begin{enumerate}
		\item $\varphi$ является гомоморфизмом аддитивных групп;
		\item $\forall a, b \in R \colon \varphi(ab) = \varphi(a)\varphi(b)$;
		\item $\varphi(1_R) = 1_S$.
	\end{enumerate}
\end{Definition}

\begin{Example}{Примеры гомоморфизмов}{}
	\begin{enumerate}
		\item Если $S \subset R$ --- подкольцо, то отображение включения $i_S\colon S \to R$, заданное как $a \mapsto a$, является гомоморфизмом.
		\item Если $I \subset R$ --- идеал, то каноническая проекция $\pi_I\colon R \to R/I$, заданная как $a \mapsto \bar{a}$, является гомоморфизмом.
	\end{enumerate}
\end{Example}

\begin{Remark}{}{}
	Аналогично теории групп, для любого гомоморфизма колец $\varphi$ можно построить коммутативную диаграмму, разложив его в композицию:
	\[ \varphi = i_{\operatorname{Im}\varphi} \circ \tilde{\varphi} \circ \pi_{\ker\varphi} \]
\end{Remark}

\begin{Theorem}{О гомоморфизме колец}{}
	Пусть $\varphi\colon R \to S$ --- гомоморфизм колец. Тогда, индуцированный гомоморфизм групп $\tilde{\varphi}\colon R/\ker\varphi \to \operatorname{Im}\varphi$ является изоморфизмом колец. При этом $\ker\varphi$ является идеалом в $R$, а $\operatorname{Im}\varphi$ --- подкольцом в $S$.
\end{Theorem}

\begin{proof}
	Из теории групп известно, что $\tilde{\varphi}$ корректно определено и является изоморфизмом аддитивных групп. Остается проверить свойства, связанные с умножением и единицей.
	
	\begin{enumerate}
		\item Покажем, что $\ker\varphi$ --- идеал в $R$. Пусть $x \in R$, $a \in \ker\varphi$. Тогда:
		\[ \varphi(xa) = \varphi(x)\varphi(a) = \varphi(x) \cdot 0 = 0 \]
		Следовательно, $xa \in \ker\varphi$.
		
		\item Покажем, что $\operatorname{Im}\varphi$ --- подкольцо в $S$. Проверим замкнутость относительно умножения и наличие единицы:
		\begin{align*}
			\varphi(a)\varphi(b) &= \varphi(ab) \in \operatorname{Im}\varphi \\
			1_{\operatorname{Im}\varphi} &= 1_S = \varphi(1_R) \in \operatorname{Im}\varphi
		\end{align*}
		
		\item Проверим сохранение умножения для $\tilde{\varphi}$:
		\[ \tilde{\varphi}(\bar{a} \cdot \bar{b}) = \tilde{\varphi}(\overline{ab}) = \varphi(ab) = \varphi(a)\varphi(b) = \tilde{\varphi}(\bar{a})\tilde{\varphi}(\bar{b}) \]
		
		\item Проверим сохранение единицы:
		\[ \tilde{\varphi}(1_{R/\ker\varphi}) = \tilde{\varphi}(\bar{1}_R) = \varphi(1_R) = 1_S = 1_{\operatorname{Im}\varphi} \]
	\end{enumerate}
	Таким образом, $\tilde{\varphi}$ является изоморфизмом колец.
\end{proof}

\subsection{Факторкольца области главных идеалов}

\begin{Remark}{}{}
	В любом кольце $R$ равенство $1 = 0$ равносильно тому, что $R = \{0\}$.
	Действительно:
	\[ 1 = 0 \iff \forall a \in R \colon a \cdot 1 = a \cdot 0 \iff \forall a \in R \colon a = 0 \iff R = \{0\} \]
\end{Remark}

\begin{Proposition}{}{}
	Пусть $R$ --- область главных идеалов (ОГИ), $b \in R$. Тогда факторкольцо $R/(b)$ является полем тогда и только тогда, когда элемент $b$ неприводим.
\end{Proposition}

\begin{proof}
	Рассмотрим различные случаи для элемента $b$.
	
	Если $b$ не является неприводимым, возможны два варианта:
	\begin{enumerate}
		\item[$1^\circ$)] $b \in R^*$ (обратим). Тогда $1 \in (b)$, откуда $(b) = R$. Следовательно, $R/(b) = \{0\}$, что по определению не является полем.
		\item[$2^\circ$)] $b = b_1 b_2$, где $b_1, b_2 \notin R^*$. Предположим от противного, что $R/(b)$ --- поле. В факторкольце имеем:
		\[ \bar{b}_1 \cdot \bar{b}_2 = \bar{b} = \bar{0} \]
		Так как в поле нет делителей нуля, то либо $\bar{b}_1 = \bar{0}$, либо $\bar{b}_2 = \bar{0}$. Пусть, без ограничения общности, $\bar{b}_1 = \bar{0}$. Это означает, что $b \mid b_1$, то есть существует $c \in R$, для которого $b_1 = bc$. Тогда:
		\[ b = b_1 b_2 = (bc) b_2 = b (cb_2) \]
		Так как $b \neq 0$, можно сократить на $b$, получив $1 = cb_2$. Это означает, что $b_2 \in R^*$, что противоречит предположению.
	\end{enumerate}
	Следовательно, если $b$ не неприводим, $R/(b)$ не может быть полем.
	
	Пусть теперь $b$ неприводим ($3^\circ$). 
	Рассмотрим ненулевой класс $\bar{c} \in R/(b)$, $\bar{c} \neq \bar{0}$. Это равносильно тому, что $c \notin (b)$. Так как $b$ неприводим и не делит $c$, их наибольший общий делитель равен единице: $(c, b) = 1$ в кольце $R$.
	
	Поскольку $R$ --- ОГИ, по расширенному алгоритму Евклида существуют элементы $x, y \in R$ такие, что:
	\[ cx + by = 1 \]
	Перейдем к классам смежности в $R/(b)$:
	\[ \bar{c} \cdot \bar{x} + \bar{b} \cdot \bar{y} = \bar{1} \]
	Так как $\bar{b} = \bar{0}$, получаем $\bar{c} \cdot \bar{x} = \bar{1}$. Следовательно, класс $\bar{c}$ обратим ($\bar{c} \in (R/(b))^*$).
	
	Таким образом, любой ненулевой элемент в $R/(b)$ обратим, значит, $R/(b)$ --- поле.
\end{proof}

\begin{Example}{Примеры факторколец}{}
	\begin{enumerate}
		\item Пусть $R = \mathbb{Z}$. Элемент $p \in \mathbb{Z}$ является неприводимым тогда и только тогда, когда $p$ --- простое число. Факторкольцо $\mathbb{F}_p := \mathbb{Z}/(p)$ является полем из $p$ элементов.
		
		\item Пусть $R = K[x]$, где $K$ --- поле. Если $f$ --- неприводимый многочлен, то факторкольцо $K[x]/(f)$ является полем. Рассмотрим частные случаи:
		
		\begin{enumerate}
			\item[(a)] $\mathbb{R}[x]/(x^2+1)$ --- поле. Классы смежности имеют вид $\{a\bar{x} + b \mid a, b \in \mathbb{R}\}$. 
			Операции задаются следующим образом:
			\begin{align*}
				(a_1\bar{x} + b_1) + (a_2\bar{x} + b_2) &= (a_1 + a_2)\bar{x} + (b_1 + b_2) \\
				\bar{x} \cdot \bar{x} &= \bar{x}^2 = \overline{-1} = -1
			\end{align*}
			
			\item[(b)] Поле $K[x]/(x-a)$. Рассмотрим гомоморфизм вычисления значения $\varphi_a\colon K[x] \to K$, заданный правилом $f \mapsto f(a)$.
			Поскольку для любой константы $c \in K$ выполнено $\varphi_a(c) = c$, гомоморфизм сюръективен, то есть $\operatorname{Im}\varphi_a = K$. Ядро гомоморфизма $\ker\varphi_a = (x-a)$. По теореме о гомоморфизме:
			\[ K[x]/(x-a) \cong K \]
			
			\item[(c)] Поле $\mathbb{R}[x]/(f)$, где $f$ --- неприводимый многочлен степени 2. Пусть $a \in \mathbb{C}$ --- корень многочлена $f$. 
			Рассмотрим гомоморфизм $\varphi\colon \mathbb{R}[x] \to \mathbb{C}$, $f \mapsto f(a)$. 
			Так как $\operatorname{Im}\varphi$ содержит $\mathbb{R}$ и содержит недействительный корень $a$, получаем $\operatorname{Im}\varphi = \mathbb{C}$. 
			Ядро $\ker\varphi = (f_0)$, где $f_0$ --- неприводимый многочлен (так как $\operatorname{Im}\varphi$ --- поле). 
			Поскольку $f \in \ker\varphi$, имеем $f_0 \mid f$, откуда $f_0 \sim f$. Значит, $\ker\varphi = (f)$, и по теореме о гомоморфизме:
			\[ \mathbb{R}[x]/(f) \cong \mathbb{C} \]
			
			\item[(d)] В кольце $\mathbb{Q}[x]$ существуют неприводимые многочлены любой степени, например $x^n - 2$.
			Факторкольцо $\mathbb{Q}[x]/(x^n - 2)$ является полем, состоящим из элементов вида:
			\[ \{a_0 + a_1\bar{x} + \dots + a_{n-1}\bar{x}^{n-1} \mid a_i \in \mathbb{Q}\} \]
			Это поле является $n$-мерным линейным пространством над $\mathbb{Q}$. Для различных значений $n$ эти поля имеют разную размерность и, следовательно, не изоморфны друг другу.
			
			\item[(e)] Пусть $K = \mathbb{F}_2$. Найдем неприводимый многочлен минимальной степени. Многочлены степени 2 в $\mathbb{F}_2[x]$:
			\begin{align*}
				x^2 &= x \cdot x \\
				x^2 + 1 &= (x + 1)^2 \\
				x^2 + x &= x(x + 1) \\
				x^2 + x + 1 & \quad \text{--- неприводим (нет корней в } \mathbb{F}_2 \text{)}
			\end{align*}
			Следовательно, $K' = \mathbb{F}_2[x]/(x^2 + x + 1)$ --- поле. Мощность этого поля $|K'| = 4$, оно состоит из элементов $\{\bar{0}, \bar{1}, \bar{x}, \overline{x+1}\}$.
		\end{enumerate}
	\end{enumerate}
\end{Example}

\begin{Remark}{}{}
	Для факторкольца $K[x]/(f)$ имеет место вложение исходного поля $K$:
	\[ K[x]/(f) \supset \{\bar{c} \mid c \in K\} = \bar{K} \cong K \]
\end{Remark}

\subsection{Подполя и характеристика поля}

\begin{Definition}{Подполе и простое поле}{}
	Подмножество $F \subset K$ называется подполем, если $F$ является подкольцом в $K$ и само является полем.
	Поле $K$ называется простым, если оно не содержит собственных подполей.
\end{Definition}

\begin{Definition}{Характеристика кольца}{}
	Пусть $A$ --- кольцо. Характеристикой кольца $A$ называется число $n$, определяемое как:
	\[ \operatorname{char} A = \min \{n \in \mathbb{N} \mid \underbrace{1 + \dots + 1}_{n \text{ раз}} = 0\} \]
	Если такого $n$ не существует, полагают $\operatorname{char} A = 0$.
\end{Definition}

\begin{Lemma}{О характеристике поля}{}
	Если $K$ --- поле, то $\operatorname{char} K = 0$ или $\operatorname{char} K = p$, где $p$ --- простое число.
\end{Lemma}

\begin{proof}
	Пусть $\operatorname{char} K = a \cdot b$, где $a > 1, b > 1$. По определению характеристики:
	\[ \underbrace{(1 + \dots + 1)}_{a \text{ раз}} \cdot \underbrace{(1 + \dots + 1)}_{b \text{ раз}} = 0 \]
	Так как в поле нет делителей нуля, хотя бы одна из скобок равна нулю. Это противоречит минимальности характеристики $\operatorname{char} K = ab$. Следовательно, характеристика либо равна нулю, либо является простым числом.
\end{proof}

\begin{Theorem}{О классификации простых полей}{}
	Пусть $K$ --- простое поле. Тогда:
	\begin{enumerate}
		\item Если $\operatorname{char} K = p \neq 0$, то $K \cong \mathbb{F}_p$.
		\item Если $\operatorname{char} K = 0$, то $K \cong \mathbb{Q}$.
	\end{enumerate}
\end{Theorem}

\begin{proof}
	Рассмотрим отображение $\varphi\colon \mathbb{Z} \to K$, заданное правилом:
	\[
	\varphi(n) = 
	\begin{cases} 
		\underbrace{1 + \dots + 1}_{n \text{ раз}}, & n > 0 \\
		0, & n = 0 \\
		-(\underbrace{1 + \dots + 1}_{-n \text{ раз}}), & n < 0 
	\end{cases}
	\]
	Отображение $\varphi$ является гомоморфизмом колец. Ядро этого гомоморфизма зависит от характеристики поля $K$:
	\[
	\ker\varphi = 
	\begin{cases} 
		(p), & \text{если } \operatorname{char} K = p \neq 0 \\
		0, & \text{если } \operatorname{char} K = 0 
	\end{cases}
	\]
	
	Рассмотрим два случая:
	\begin{enumerate}
		\item Пусть $\operatorname{char} K = p \neq 0$. По теореме о гомоморфизме:
		\[ \mathbb{Z}/(p) \cong \operatorname{Im}\varphi \]
		Поскольку $\mathbb{Z}/(p)$ --- поле, $\operatorname{Im}\varphi$ также является полем и, следовательно, подполем в $K$. Так как $K$ --- простое поле (не содержит собственных подполей), получаем $\operatorname{Im}\varphi = K$. Таким образом, $K \cong \mathbb{Z}/(p) = \mathbb{F}_p$.
		
		\item Пусть $\operatorname{char} K = 0$. Тогда $\ker\varphi = 0$, следовательно, $\varphi$ --- инъективно.
		Построим отображение $\varphi_1\colon \mathbb{Q} \to K$, заданное правилом $\varphi_1\left(\frac{a}{b}\right) = \frac{\varphi(a)}{\varphi(b)}$. (Так как $b \neq 0$, то $\varphi(b) \neq 0$).
		Проверим корректность:
		\[ \frac{a}{b} = \frac{a'}{b'} \implies ab' = a'b \implies \varphi(a)\varphi(b') = \varphi(a')\varphi(b) \implies \frac{\varphi(a)}{\varphi(b)} = \frac{\varphi(a')}{\varphi(b')} \]
		Легко проверить, что $\varphi_1$ является гомоморфизмом колец. 
		Найдем ядро: $\varphi_1\left(\frac{a}{b}\right) = 0 \implies \varphi(a) = 0 \implies a = 0 \implies \frac{a}{b} = 0$. Значит, $\ker\varphi_1 = 0$.
		По теореме о гомоморфизме $\mathbb{Q} \cong \operatorname{Im}\varphi_1$. 
		Следовательно, $\operatorname{Im}\varphi_1$ --- подполе в $K$. Так как $K$ --- простое поле, $\operatorname{Im}\varphi_1 = K$, откуда $K \cong \mathbb{Q}$.
	\end{enumerate}
\end{proof}

\begin{Proposition}{}{}
	Любое поле содержит простое подполе.
\end{Proposition}

\begin{proof}
	Пусть $K$ --- поле. Рассмотрим пересечение всех его подполей:
	\[ F_0 := \bigcap_{F \text{ --- подполе } K} F \]
	Очевидно, что $F_0$ само является подполем в $K$. Любое подполе $F \subset K$ содержит $F_0$, поэтому $F_0$ не содержит собственных подполей и по определению является простым полем.
\end{proof}